\documentclass[11pt,a4paper]{article}

% --- Encoding & fonts ---
\usepackage[T1]{fontenc}
\usepackage[utf8]{inputenc}
\usepackage{mathpazo}              % Palatino for text and math
\usepackage[scaled=0.85]{helvet}   % Helvetica for sans-serif
\usepackage{microtype}             % Improved typographic refinement

% --- Mathematics ---
\usepackage{amsmath,amssymb,amsfonts}
\usepackage{slantsc}

% --- Colors & page layout ---
\usepackage[dvipsnames]{xcolor}
\usepackage[margin=25mm, top=30mm, bottom=30mm]{geometry}

% --- Headers & footers ---
\usepackage{fancyhdr}
\pagestyle{fancy}
\fancyhf{}
\fancyhead[L]{\small\itshape Nestor (1996) --- Thesis Extracts}
\fancyhead[R]{\small\itshape Sections 2.1.2, 3.1 \& 3.2}
\fancyfoot[C]{\thepage}
\renewcommand{\headrulewidth}{0.4pt}
\renewcommand{\footrulewidth}{0pt}

% --- Section formatting ---
\usepackage{titlesec}
\titleformat{\section}{\large\bfseries}{}{0em}{}[\vspace{0.3em}\hrule\vspace{0.5em}]
\titleformat{\subsection}{\normalsize\bfseries}{}{0em}{}
\titleformat{\subsubsection}{\normalsize\itshape\bfseries}{}{0em}{}

% --- Captions ---
\usepackage[font=small, labelfont=bf, labelsep=period]{caption}

% --- Spacing ---
\usepackage{setspace}
\setstretch{1.15}
\setlength{\parskip}{0.4em}

% --- Hyperlinks (must be loaded last, before macros) ---
\usepackage[colorlinks=true, linkcolor=blue!60!black, citecolor=blue!60!black, urlcolor=blue!60!black]{hyperref}

% --- Macros from the thesis ---
\newcommand{\CB}[1]{\mbox{\boldmath$\cal #1$\unboldmath}}
\newcommand{\bmi}[1]{\mbox{\boldmath$\mit #1$\unboldmath}}
\newcommand{\rf}[1]{\stackrel{\scriptscriptstyle _{\,\,\circ}}{ #1}}
\newcommand{\gtlt}[0]{\stackrel{\scriptscriptstyle >}{\scriptscriptstyle _<}}

\begin{document}

% --- Title block ---
\thispagestyle{plain}
\begin{center}
\vspace*{1em}
{\LARGE\bfseries Governing Equations and Representations of\\[0.3em]
the Free Space Green's Tensor}\\[1.8em]
{\large T.\,M.\ Nestor}\\[0.5em]
{\itshape Research School of Earth Sciences, Australian National University}\\[2em]
\end{center}

\noindent\fbox{\parbox{\dimexpr\textwidth-2\fboxsep-2\fboxrule}{%
\small\textbf{Source:}\enspace Extracted from Chapters 2 and 3 of:
\textit{Seismic Wave Propagation in Complex Media: Summing the
Multiple Scattering Series}, PhD Thesis, ANU (1996).
Sections 2.1.2, 3.1, and 3.2 are reproduced here.}}


\section*{2.1.2\quad The Governing Equations}


Assuming that the depth below the free surface is measured by the
$z-$coordinate, the forced elastic displacement ${\bmi  u} = {({u}_{z},{u}_{x},{u}_{y})}^{T}$
and traction vector ${\bmi  t} = {({\tau}_{zz},{\tau}_{xz},{\tau}_{yz})}^{T }$ on a
fixed depth plane (i.e. $z=const.$) with unit
normal ${\bmi  e}_{z}$,  are grouped into the $6$-vector
\begin{equation} \label{bdef}
{\bmi  b}(z,x,y)= {\left[ \begin{array}{ccc}{ u}_{ z}
(z,x,y),&{u}_{x}(z,x,y),&{u}_{y}(z,x,y),\\
{\tau}_{zz}(z,x,y),&{\tau}_{xz}(z,x,y),&
{\tau}_{yz}(z,x,y)\end{array} \right]}^{  T }.
\end{equation}
The governing equation for time harmonic elastic wave motion in an {\em isotropic  medium\/}
is then a
set of $6$ coupled partial differential equations (first order with respect
to the depth coordinate) which may be written in the operator notation
\begin{equation} \label{govEqn00}
{\CB L}({\bmi x};{\partial }_{z},{\partial }_{x},{\partial }_{y})\cdot {\bmi  b }\equiv\left[\partial_{z}-{\CB A}({\bmi x};{\partial }_{x},{\partial }_{y})\right]\cdot {\bmi  b }=
{\bmi  f}
\end{equation}
where ${\bmi  b}$ is given in (\ref{bdef})  and
${\CB A}$ is defined as
\begin{equation} \label{Adef0}
{\CB A}({\bmi x};{\partial }_{x},{\partial }_{y})=\left[{\begin{array}{cccccc}
0&- \gamma {\partial }_{  x}&  - \gamma {\partial }_{  y}&  a&0&0\\
-{ \partial }_{x}&0&0&0&b&0\\
-{\partial }_{y}&0&0&0&0&b\\
- \rho {\omega }^{  2}&  0&0&0&-{ \partial }_{x}&-{ \partial }_{y}\\
0&- \rho {\omega }^{  2}  -{\partial }_{x} \zeta {\partial }_{  x}  -{ \partial }_{y} \mu
{\partial }_{  y}&  -{ \partial }_{x} \chi {\partial }_{  y}
 -{ \partial }_{y} \mu {\partial }_{  x}&  -{ \partial }_{x} \gamma &  0&0\\
 0&-{ \partial }_{y} \chi {\partial }_{  x}  -{ \partial }_{x} \mu {\partial }_{  y}&
   - \rho {\omega }^{  2}  -{ \partial }_{y} \zeta {\partial }_{  y}
 -{ \partial }_{x} \mu {\partial }_{  x}&  -{ \partial }_{y} \gamma &  0&0\end{array}}
 \right]
\end{equation}
and the material parameters in (\ref{Adef0}) are defined in Table \ref{ATable}.

\begin{table}[h]
\begin{tabular*}{85mm}{ccc}\hline
parameter & $(\rho, \alpha, \beta)$-definition & $(\rho, \lambda, \mu)$-definition\\ \hline
$P$&$\alpha$&$\sqrt{{{\lambda + 2 \mu} \over {\rho}}}$\\
$S$&$\beta$&$\sqrt{{{\mu} \over {\rho}}}$\\
$H$&$\beta$&$\sqrt{{{\mu} \over {\rho}}}$\\
$\gamma$&$\left(1- 2{{\beta^2}\over{\alpha^2}}\right)$&${{\lambda}\over{\lambda+2 \mu}}$\\
$\zeta$&$4 \rho \beta^2 \left(1- {{\beta^2}\over{\alpha^2}} \right)$&${{4
\mu (\lambda+ \mu)}\over{\lambda+2 \mu}}$\\
$\chi$&$2 \rho \beta^2\left(1-
2{{\beta^2}\over{\alpha^2}}\right)$&${{2 \lambda \mu}\over{\lambda+2
\mu}}$\\ $a$&${{1}\over{\rho \alpha^2}}$&${{1}\over{\lambda+2 \mu}}$\\
$b$&${{1}\over{\rho \beta^2}}$&${{1}\over{\mu}}$\\
\hline
\end{tabular*}
\caption{\label{ATable} { The material parameter definitions}}
\end{table}

This thesis is concerned with finding solutions of (\ref{govEqn00}) when
the material parameters in Table \ref{ATable} are functions of $(z,x)$ but are
independent of  the $y-$coordinate. This problem is sometimes called the
$2\,{1/2}-$D problem because the wavefield may spread in $3-$D,
but the heterogeneity is confined to $2-$D.
We exploit the translational invariance of the medium along the
$y-$coordinate by assuming that the displacement/traction field ${\bmi
b}(z,x,y)$ may be {\em synthesized\/} from its  $k_y$ plane wave components as
\begin{equation} \label{kysyn}
{\bmi b}(z,x,y)={\frac{1}{2\pi }}\int_{R} \, {\rm d}{k}_{y}\,{\bmi
b}(z,x,{k}_{y})\, {\rm e}^{i{k}_{y}y}
\end{equation}
where ${\bmi b}(z,x,{k}_{y})$ is now governed by
\begin{equation} \label{govEqn0}
{\CB L}(z,x;{\partial }_{z},{\partial }_{x},i {k}_{y})\cdot {\bmi  b
}(z,x,k_y; \omega)\equiv
\left[\partial_{z}-{\CB A}(z,x;{\partial }_{x},i {k}_{y})\right]\cdot {\bmi  b }(z,x,k_y;
\omega)=
\tilde{\bmi  f}
\end{equation}
with
\begin{equation} \label{Adef}
{\CB A}(z,x;{\partial }_{x},i {k}_{y})=\left[{\begin{array}{cccccc}
0&- \gamma {\partial }_{  x}&  -i \gamma {  k}_{  y}&  a&0&0\\ -{ \partial
}_{x}&0&0&0&b&0\\ -i{k}_{y}&0&0&0&0&b\\ - \rho {\omega }^{
2}&  0&0&0&-{ \partial }_{x}&-i{k}_{y}\\ 0&- \rho {\omega
}^{  2}  -{ \partial }_{x} \zeta {\partial }_{
x}  + \mu {  k}_{  y}^{  2}&
-i{k}_{y}{ \partial }_{x} \chi   -i{k}_{y} \mu {\partial
}_{  x}&  -{ \partial }_{x} \gamma &  0&0\\
0&-i{k}_{y} \chi {\partial }_{  x}  -i{k}_{y}{ \partial
}_{x} \mu &  - \rho {\omega }^{  2}  + \zeta
{  k}_{  y}^{  2}  -{ \partial }_{x} \mu
{\partial }_{  x}&  -i{k}_{y} \gamma &  0&0\end{array}}\right]
\end{equation}
and
\begin{equation} \label{kysynf}
\tilde{\bmi f}=\int_{R} \, {\rm d}{k}_{y}\,{\bmi f}(z,x,y)
\, {\rm e}^{-i{k}_{y}y}.
\end{equation}
In this work, we assume the the motion is driven by a point source in the
$y=0$ plane, i.e. ${\bmi f}(z,x,y)= {\bmi F}(z,x) \delta(y)$, so that
from (\ref{kysynf}) we have
\begin{equation} \label{kysynf1}
\tilde{\bmi f}=\int_{R} \, {\rm d}{k}_{y}\,{\bmi F}(z,x) \delta(y)
\, {\rm e}^{-i{k}_{y}y} = {\bmi F}(z,x).
\end{equation}
Due to the translational invariance of the medium in the $y-$coordinate,
if the source plane is translated to  $y=y_1$, then the spatial domain solution
${\bmi  b }(z,x,y; \omega)$ is translated to  ${\bmi  b }(z,x,y-y_1;
\omega)$. Consequently, it is straight forward to derive the solution
appropriate to general source conditions, by translating the field
for a point source in the $y=0$ plane, accordingly.

In (\ref{govEqn0}) both $\omega$ and ${  k}_{  y}$ are {\em passive
parameters\/}. The space/time solution is constructed from a
superposition of solutions with different values of $\omega$ and ${  k}_{  y}$
as prescribed by the {\em inverse Fourier transforms\/}
on frequency and on $k_y$. The main task remaining, is the construction the
transformed displacement/traction vector ${\bmi  b }(z,x,k_y; \omega)$ in (\ref{govEqn0}).
Often, we will not indicate the explicit dependence of  ${\bmi  b }(z,x,k_y; \omega)$
on $k_y$ and $\omega$ but write ${\bmi  b }(z,x)$ instead.

The representation (\ref{govEqn0}) is most convenient when the
dominant material variation is in the $z$ coordinate direction, as we have
singled out this coordinate for special treatment, allowing easy application of
boundary and interface conditions on fixed depth planes.  Also, the
displacement ${\bmi  u}$ and depth plane traction ${\bmi  t}$ are accorded equal status.
In the more traditional displacement formulation, the  offset $x$ and depth $z$
coordinates appear within the equations in a symmetric form (due to isotropy)
whereas the traction field must be derived from the displacement via the
elastic analogue of {\em Hooke's law\/} which involves further spatial
differentiations of $\bmi  u$.


We note in passing that (\ref{govEqn0}) simplifies considerably in the
special case where the motion is driven by a {\em line source\/}, along
the $y-$axis so that we may set  $k_y\equiv 0$ in (\ref{Adef}) which then
reduces to
\[
{\CB A}(z,x;{\partial }_{x},0)
=\left[{\begin{array}{cccccc}
0&- \gamma {\partial }_{  x}&  0&  a&0&0\\ -{ \partial
}_{x}&0&0&0&b&0\\ 0&0&0&0&0&b\\ - \rho {\omega }^{
2}&  0&0&0&-{ \partial }_{x}&0\\ 0&- \rho {\omega
}^{  2}  -{ \partial }_{x} \zeta {\partial }_{
x} &
0&  -{ \partial }_{x} \gamma &  0&0\\
0&0 &  - \rho {\omega }^{  2}
 -{ \partial }_{x} \mu
{\partial }_{  x}& 0 &  0&0\end{array}}\right].
\]
In this instance,  (\ref{govEqn0}) may be decomposed into its {\em in plane\/}
equation
\begin{equation} 	\label{PSVeqn}
{\frac{ \partial }{\partial z}}\left[{\begin{array}{c}{u}_{z}\\
{u}_{x}\\
{\tau }_{zz}\\
{\tau }_{xz}\end{array}}\right]
=\left[{\begin{array}{cccc}
0&- \gamma {\partial }_{  x}&  a&0\\
-{ \partial }_{x}&0&0&b\\
- \rho {\omega }^{  2}&  0&0&-{ \partial }_{x}\\
0&- \rho {\omega }^{  2}  -{ \partial }_{x} \zeta {\partial }_{  x} &
 -{ \partial }_{x} \gamma &  0
\end{array}}\right]
\left[{\begin{array}{c}{u}_{z}\\
{u}_{x}\\
{\tau }_{zz}\\
{\tau }_{xz}\end{array}}\right]
\end{equation}
for the $P/SV$ field and {\em out of plane\/} equation
\begin{equation} 	\label{SHeqn}
	{\frac{\partial }{\partial z}}\left[{\begin{array}{c}{u}_{y}\\
	{\rm \tau }_{yz}\end{array}}\right]
	=\left[{\begin{array}{cc}
	0&b\\
	- \rho {\omega }^{  2}-{ \partial }_{x} \mu {\partial }_{  x}& 0
	\end{array}}\right]
	\left[{\begin{array}{c}{u}_{y}\\
	{\rm \tau }_{yz}\end{array}}\right]
\end{equation}
for the $SH$ field.
For the most part, we will concentrate on the solution of (\ref{govEqn0})
because the source conditions needed for the applicability of (\ref{PSVeqn}) and (\ref{SHeqn})
are too restrictive.

Consistent with the decomposition of the material properties, we
 {\em split\/} (\ref{Adef}) into
\begin{equation} \label{Asplit}
{\CB A}({\bmi x};{\partial }_{x},{\partial }_{y};k_y)=
{\rf{\bmi A}}({\bmi x};{\partial }_{x},{\partial }_{y})
+{{\mit \Delta} {\bmi A}}({\bmi x};{\partial }_{x},{\partial }_{y})
\end{equation}
where
\begin{equation} \label{Adefref}
{\rf{\bmi A}}(z,x;{\partial }_{x},{k}_{y})=\left[{\begin{array}{cccccc}
0&- {\rf \gamma} {\partial }_{  x}&  -i {\rf \gamma} {  k}_{  y}&  {\rf a}&0&0\\ -{ \partial
}_{x}&0&0&0&{\rf b}&0\\ -i{k}_{y}&0&0&0&0&{\rf b}\\ - {\rf \rho} {\omega }^{
2}&  0&0&0&-{ \partial }_{x}&-i{k}_{y}\\ 0&- {\rf \rho} {\omega
}^{  2}  -{ \partial }_{x} {\rf \zeta} {\partial }_{
x}  + {\rf \mu} {  k}_{  y}^{  2}&
-i{k}_{y}{ \partial }_{x} {\rf \chi}   -i{k}_{y} {\rf \mu} {\partial
}_{  x}&  -{ \partial }_{x} {\rf \gamma} &  0&0\\
0&-i{k}_{y} {\rf \chi} {\partial }_{  x}  -i{k}_{y}{ \partial
}_{x} {\rf \mu} &  - {\rf \rho} {\omega }^{  2}  + {\rf \zeta }
{  k}_{  y}^{  2}  -{ \partial }_{x} {\rf \mu}
{\partial }_{  x}&  -i{k}_{y} {\rf \gamma} &
0&0\end{array}}\right]
\end{equation}
where ${\rf \gamma}$ is the {\em reference medium\/} value of ${\gamma}$
from Table \ref{ATable} and similar definitions hold for the remaining
{\em reference\/} material parameters
in (\ref{Adefref}). Also,
\begin{equation} \label{dAdef}
{{\mit \Delta} {\bmi A}}(z,x;{\partial }_{x},{k}_{y})=\left[{\begin{array}{cccccc}
0&- {{\mit \Delta} \gamma} {\partial }_{  x}&  -i {{\mit \Delta} \gamma} {  k}_{  y}&
{{\mit \Delta} a}&0&0\\0&0&0&0&{{\mit \Delta} b}&0\\ 0&0&0&0&0&{{\mit \Delta} b}\\ - {{\mit \Delta} \rho} {\omega }^{
2}&  0&0&0&0&0\\ 0&- {{\mit \Delta} \rho} {\omega
}^{  2}  -{ \partial }_{x} {{\mit \Delta} \zeta} {\partial }_{
x}  + {{\mit \Delta} \mu} {  k}_{  y}^{  2}&
-i{k}_{y}{ \partial }_{x} {{\mit \Delta} \chi}   -i{k}_{y} {{\mit \Delta} \mu} {\partial
}_{  x}&  -{ \partial }_{x} {{\mit \Delta} \gamma} &  0&0\\
0&-i{k}_{y} {{\mit \Delta} \chi} {\partial }_{  x}  -i{k}_{y}{ \partial
}_{x} {{\mit \Delta} \mu} &  - {{\mit \Delta} \rho} {\omega }^{  2}  + {{\mit \Delta} \zeta }
{  k}_{  y}^{  2}  -{ \partial }_{x} {{\mit \Delta} \mu}
{\partial }_{  x}&  -i{k}_{y} {{\mit \Delta} \gamma} &  0&0\end{array}}\right]
\end{equation}
with the material contrast parameters in (\ref{dAdef}) defined in Table \ref{contrastTable}.

\begin{table}[h]
\begin{tabular*}{100mm}{ll}\hline
contrast & definition\\ \hline
 $ \Delta \rho$ & $\rho - {\rf \rho}$\\
 $ \Delta \mu$ & $\mu - {\rf \mu}$\\
 $ \Delta \gamma$ & $2{\frac{{ \rf \beta }{}^{2}}{{ \rf \alpha }{}^{2}}}-
 2{\frac{{\beta}^{2}}{{\alpha}^{2}}}$\\
$\Delta \zeta$ &$4 \rho{\beta
}^{2}\left[{1-{\frac{{\beta  }^{2}}{{\alpha}^{2}}}}\right]-4{\rf \rho}
{\rf \beta }{}^{  2}\left[{
 1-{\frac{{ \rf \beta }{}^{2}}{{\rf  \alpha }{}^{2}}}}\right]$\\
$\Delta \chi$ &$2 \rho {\beta}^{2}\left[{1-2{\frac{{ \beta
}^{2}}{{\alpha}^{2}}}}\right]-2{\rf \rho} {\rf \beta }{}^{  2}\left[{
 1-2{\frac{{\rf \beta }{}^{2}}{{\rf \alpha }{}^{2}}}}\right]$\\
$\Delta  a$ & ${\frac{1}{\rho{\alpha}^{2}}}-
 {\frac{1}{ {\rf \rho} {{\rf \alpha} }{}^{  2}}}$\\
$\Delta  b$ & ${\frac{1}{\rho{\beta}^{2}}}-
 {\frac{1}{ {\rf \rho} {{\rf \beta} }{}^{  2}}}$\\
 \hline
\end{tabular*}
\caption{\label{contrastTable} { The material contrast parameters}}
\end{table}

Using the decomposition of (\ref{Asplit}) in the governing equation (\ref{govEqn0}) gives
\begin{equation} \label{Leqn}
{\bmi L}(z,x;{\partial }_{z},{\partial }_{x},{k}_{y})\cdot{\bmi b}(z,x,k_y; \omega) =
\tilde{\bmi f} + {\mit \Delta} {\bmi A}(z,x;{\partial }_{x},{k}_{y})\cdot{\bmi b}(z,x,k_y; \omega)
\end{equation}
where the {\em reference medium\/} operator ${\bmi L}$ is defined as
\begin{equation} \label{Lref}
{\bmi L}(z,x;{\partial }_{z},{\partial }_{x},{k}_{y})\equiv
\left[\partial_{z}-\rf{\bmi A}(z,x;{\partial }_{x},{k}_{y})\right]
\end{equation}
with $\rf{\bmi A}$ given in (\ref{Adefref}).
Scrupulously indicating all the arguments of the various operators introduced in
this thesis can be very cumbersome, so for example, we will often write
${\bmi L}$ to mean ${\bmi
L}(z,x;{\partial }_{z},{\partial }_{x},{k}_{y})$ when no misunderstanding
is likely.

%%%%%%%%%%%%%%%%%%%%%%%%%%%%%%%%%%%%%%%%%%%%%%%%%%%%%%%%%%%%%%%%%%%%%
\newpage

\section*{Chapter 3: Representations of the Time Harmonic Elastic Wave Green's Tensor}

We recall from Chapter 2 that
 the material
properties of the propagation environment have been defined as
the sum of a reference and residual component, i.e.
$(\rho, \alpha, \beta)=(\rf{\rho}+{\mit \Delta}\rho,\rf{\alpha}+{\mit
\Delta}\alpha,\rf{\beta}+{\mit \Delta}\beta)$. In this chapter, we will
be concerned exclusively with uniform media so that
residual properties $({\mit \Delta}\rho,{\mit
\Delta}\alpha,{\mit \Delta}\beta)$ vanish and we have
$(\rho, \alpha, \beta)=(\rf{\rho},\rf{\alpha},\rf{\beta})$. To simplify
notation, we discontinue the practice of writing  $(\rf{\rho},\rf{\alpha},\rf{\beta})$
to denote reference medium parameters and write
$({\rho},{\alpha},{\beta})$ instead.

This chapter is concerned with the construction of the time harmonic
elastic wave motion driven by a {\em lattice of multipolar point forces\/}
${\bmi f}$ {\em  superimposed on a homogeneous medium\/}.
The homogeneous medium is
characterized by the  elastic parameters $(\rho,\alpha,\beta)$, which
denote the density, compressional and shear speeds respectively, in an
isotropic medium.

Each point force  ${\bmi f}$,  has an assumed time
dependence of the form
$\exp(\epsilon t -i 2 \pi \nu t)$ where
$\nu $ is the  frequency and $\epsilon$ is a non-negative {\em
growth\/} parameter. The {\em adiabatic growth term\/} ${\rm e}^{\epsilon
t}$,
evolves the wavefield from a state of quiescence in the infinite past, to
its
present value, following a path that ensures that the {\em wavefield is
always in steady
state equilibrium with the driving force\/} ${\bmi f}$. This device has been
employed by
{\em Lighthill} ~(1960) to enforce the {\em radiation
boundary conditions\/}.
For  convenience, we introduce the {\em complex angular frequency\/}
\begin{equation} \label{w+}
\omega\equiv 2 \pi \nu +i \epsilon.
\end{equation}
We can then synthesize {\em transient motion\/} from a discrete set of
{\em
forced motions\/}, sampled from the {\em Bromwich contour\/} in the
complex frequency
domain [e.g.\ {\em Chapman} ~(1978)] as
\begin{equation}	\label{InvLaplace}
	\left[{\cdot }\right](t)=2 \, {\rm e}^{\epsilon  t}\,
	\Re \int_{ 0}^{\nu_{max}} {\rm d}\nu \ {\rm e}^{ -i
	2 \pi \nu t}\left[{\cdot }\right] (2 \pi \nu +i \epsilon)
\end{equation}
where $\nu_{{max}}$ is the maximum frequency of the source field.
Usually the growth parameter $\epsilon$ is chosen so that $5 \le \epsilon T
\le 10$, where $T$ is the length of the required time series (which is
usually determined by trial and error). Using the complex angular frequency $\omega$ smoothes the frequency domain
response, which enhances the conditioning of the problem,
causing the complex frequency multiple scattering series to converge
faster than the real frequency version. However, we must be aware that the
process is equivalent to an analytic continuation, and is inherently
unstable. The instability, is manifested by the sensitivity of
(\ref{InvLaplace}) to errors in the frequency domain response. This is a
particularly troubling issue, as the multiple scattering series is
usually only summed to a relative error of $10^{-6}$ and it remains an
open question whether we have secured sufficient accuracy to avoid the
deleterious effects of the numerical instability inherent in
(\ref{InvLaplace}). An excellent survey of numerical methods for the
numerical evaluation of (\ref{InvLaplace}) is presented in
{\em Davies \& Martin\/}~(1979).

\subsection*{3.1\quad The Free Space, Time Harmonic, Elastic Wave Green's Tensor}
Given the 2-D vectors
\[
{\bmi  x}\equiv (z,x), \quad {\bmi  x}^{\prime}\equiv (z^{\prime},x^{\prime})
\]
a particularly useful solution $\rf{\bmi  G }({\bmi  x};{\bmi  x}^{\prime};k_y)$,
satisfies (\ref{govEqn0}) 	when the the material properties in
Table \ref{ATable} are homogeneous and the  {\em $y-$Fourier transformed body
force\/} ${\bmi F}({\bmi x};{\bmi x}^{\prime})$ assumes the
special form
\[
{\bmi F}({\bmi x};{\bmi x}^{\prime})=
{\bmi  I}_6 \delta({\bmi x}-{\bmi x}^{\prime}),
\]
with ${\bmi  I}_6$ the $6 \times 6$ {\em identity matrix\/}, {\em viz\/}
\begin{equation} \label{GreenDef}
{{\CB L}}\cdot \rf{\bmi  G}({\bmi  x};{\bmi  x}^{\prime};k_y)={\bmi  I}_6
\delta({z}-{z}^{\prime})\delta({x}-{x}^{\prime}).
\end{equation}
We can then use the linearity of (\ref{govEqn0}) to pose the {\em
particular solution\/} ${\bmi  b}_{P}$ of (\ref{govEqn0}) as the {\em
superposition integral\/}
\begin{equation}	\label{superposition}
	{\bmi  b}_{P}=\int {\rm d}{\bmi  x}^{\prime}\, \rf{\bmi  G}({\bmi  x};{\bmi
x}^{\prime};k_y){\bmi
	f}({\bmi  x}^{\prime})
\end{equation}
for an arbitrary body force $\bmi  f$.
It will be  convenient for the purpose of subsequent constructions, to
interpret
$\rf{\bmi  G}({\bmi  x};{\bmi  x}^{\prime};k_y)$  of (\ref{superposition}), as the
{\em integral kernel\/} of the formal inverse
\begin{equation}	\label{invDef}
	\left[{{\CB L}}^{-1} {\bmi  w}\right]({\bmi  x})\equiv \int {\rm d}{\bmi
	x}^{\prime}\, \rf{\bmi  G }({\bmi  x};{\bmi  x}^{\prime};k_y)\cdot{\bmi  w}({\bmi
x}^{\prime})
\end{equation}
for an arbitrary {\em test\/} vector-function $\bmi  w$, where ${{\CB L}}$ is the governing
operator in (\ref{govEqn0}). Equation
(\ref{invDef}) is strictly formal,  as the {\em null space\/} of
${{\CB L}}$ is spanned by the set of {\em homogeneous\/} solutions and
consequently,
$\rf{\bmi  G}$ is not unique. To get a {\em well posed\/} inverse, we invoke
the {\em radiation condition\/} that {\em there are no homogeneous
solutions
of\/} (\ref{govEqn0}) {\em due to sources at infinity\/}.


To solve for $\rf{\bmi  G}$ in (\ref{GreenDef}), or equivalently, to construct
the
inverse  ${{\CB L}}^{-1}$ in (\ref{invDef}), we follow the
classical {\em Fourier Transform\/} route. To begin, we transform the $x-$
dependence via ${\cal F}_{x}\, [\cdot]({k}_{x})\equiv \frac{1}{\sqrt{2
\pi}}\int_{-\infty }^{ \infty
}dx\ \exp \left({-i\, {k}_{x} x}\right)\, [\cdot](x)$, where  ${k}_{x}$ is the
horizontal
component
of the wavevector,
so that the transverse {\em Fourier\/} transform of the governing equation
(\ref{GreenDef})
may be written as
\begin{align}	\label{FLdef}
\langle { {k}_{x}} | {\CB L }| { {k}_{x}}^{\prime}\rangle
\langle { {k}_{x}}^{\prime}| x\rangle\langle x| \,\rf{\bmi
G}({\bmi x};{\bmi x}^{\prime};k_y)&\nonumber\\
&\equiv{\bmi  L}(z,{k}_{x},{k}_{y}) \,\rf{\bmi  G}(z,{k}_{x};z^{\prime},x^{\prime};k_y)\nonumber\\
&\equiv \left[\partial_z -{\bmi  A}({k}_{x},{k}_{y})\,  \right]\,\rf{\bmi
G}(z,{k}_{x};z^{\prime},x^{\prime};k_y)={\bmi  I}_{6}
	\delta(z-z^{\prime}) \, \frac{\exp(-i {k}_{x} x^{\prime})}{\sqrt{2 \pi}}
\end{align}
where the {\em Fourier domain system matrix\/} $\langle { {k}_{x}} | {\CB A }|
{ {k}_{x}}^{\prime}\rangle= \delta({ {k}_{x}}-{ {k}_{x}}^{\prime}) {\bmi
A}({k}_{x},{k}_{y})$ is defined
\begin{equation}	\label{Akdef}
	{\bmi  A}({k}_{x},{k}_{y})=\left[{\begin{array}{cccccc}
0&-i \gamma {  k}_{  x}&
 - \gamma   i\ {k}_{y}&a&0&0\\ -i{k}_{x}&0&0&0&b&0\\
-i{k}_{y}&0&0&0&0&b\\ - \rho {\omega }^{  2}&
0&0&0&-i{k}_{x}&-i{k}_{y}\\
0&- \rho {\omega }^{  2}  +
\zeta {  k}_{  x}^{  2}  + \mu {
k}_{  y}^{  2}&{  k}_{  x}{  k}_{
 y}\left({\chi   + \mu }\right)&  -i{k}_{x} \gamma
&  0&0\\ 0&{k}_{x}{k}_{y}\left({ \chi   + \mu
}\right)&- \rho {\omega }^{  2}  + \zeta {
k}_{  y}^{  2}  + \mu {  k}_{
x}^{  2}&  -i{k}_{y} \gamma &  0&0\end{array}}\right]
\end{equation}
due to the {\em translational invariance of\/} ${\CB A }$.
The system operator ${\bmi  A}({k}_{x},{k}_{y})$ of
(\ref{Akdef}) is {\em quasi-Hamiltonian\/} since
\begin{equation} \label{ATdef}
{\bmi  A}^{T}(-{k}_{x},-{k}_{y}){\bmi  J}_6 + {\bmi  J}_6 {\bmi
A}({k}_{x},{k}_{y}) \equiv {\bf 0}_6
\end{equation}
where ${\bmi  J}_6$ is the {\em canonically symplectic\/} array.
The {\em diagonal array\/} ${\bmi
\Lambda}$ {\em of eigenvalues\/}  of ${\bmi  A}({k}_{x},{k}_{y})$ is defined as
\begin{equation}	\label{eigDef}
{\bmi \Lambda}({k}_{x},{k}_{y}) = i \ \mbox{diag} \left[{k}_{z,P
	}({k}_{x},{k}_{y}),{k}_{z,S }({k}_{x},{k}_{y}),{k}_{z,H }({k}_{x},{k}_{y}),
	-{k}_{z,P }({k}_{x},{k}_{y}),-{k}_{z,S }({k}_{x},{k}_{y}),
	-{k}_{z,H }({k}_{x},{k}_{y}) \right]
\end{equation}
where
\begin{equation}	\label{kzcDef}
	{k}_{z,c}({k}_{x},{k}_{y}) = \left\{
	\begin{array}{c}
	\sqrt {\hat{ K}_{c}^{2}({k}_{y})\ -\ {k}_{x}^{2}}\ \ \ \ \ \ |{k}_{x}|\ <\
	|\hat{K}_{c}({k}_{y})|\\
	i\sqrt {{k}_{x}^{ 2}\ -\ \hat{ K}_{c}^{2}({k}_{y})} \ \ \ \ |{k}_{x}|\ \ge   \
	|\hat{K}_{c}({k}_{y})|
\end{array}
\right. , \quad  c\in\{P,S,H\},
\end{equation}
with
\begin{equation}	\label{Kdef}
	\hat{K}_{c}({k}_{y}) = \left\{
	\begin{array}{c}
	\sqrt {(\omega/{c})^{2}\ -\ {k}_{y}^{2}}\ \ \ \ \ \ |{k}_{y}|\ <\
	|\omega/{c}|\\
	i\sqrt {{k}_{y}^{ 2}\ -\ (\omega/{c})^{2}} \ \ \ \ |{k}_{y}|\ \ge   \ |\omega/{c}|
\end{array} \right. , \quad  c\in\{P,S,H\}.
\end{equation}
{\em We note that ${k}_{z,S }={k}_{z,H }$
because $S=H=\beta$ from Table \ref{ATable}. We distinguish
between ${k}_{z,S }$ and ${k}_{z,H }$,  only to maintain a self-consistent
notational convention.\/}
 The system matrix ${\bmi  A}({k}_{x},{k}_{y})$ of (\ref{Akdef}) has an equivalent {\em
spectral representation\/}:
\begin{equation}	\label{specA}
	{\bmi  A}({k}_{x},{k}_{y}) \equiv {\bmi  D}_{z}({k}_{x},{k}_{y}){\bmi
	\Lambda }({k}_{x},{k}_{y}){\bmi  D}_{z}^{ -1}({k}_{x},{k}_{y});
\end{equation}
where the {\em eigen matrix\/} ${\bmi  D}_{z}({k}_{x},{k}_{y})$ is  defined as
\begin{equation} \label{Ddef}
{\bmi D}_{z}(k_{x},k_{y})=\left[\begin{array}{cccccc}
\left|+,P, k_{x},k_{y}  \right\rangle&\left|+,S, k_{x},k_{y}
\right\rangle&\left|+,H, k_{x},k_{y}  \right\rangle&
\left|-,P, k_{x},k_{y}  \right\rangle&\left|-,S, k_{x},k_{y}
\right\rangle&\left|-,H, k_{x},k_{y}  \right\rangle
\end{array}\right]
\end{equation}
with:
\begin{itemize}
\item the $\left|\pm,P, k_{x},k_{y}  \right\rangle$ eigen-vectors
are the   ${\stackrel{\mbox{\small downgoing}}{{}_{\mbox{\small
upgoing}}}}$  $P-$waves
\begin{equation} \label{Peigen}
\left|\pm,P, k_{x},k_{y}  \right\rangle=\left[\begin{array}{c} \pm
i\,{k}_{z,P}(k_x,k_y)\\i\,k_{x}\\i\,k_{y}\\
  \rho \,\left( 2\,{\beta^2}\,{k_{x}^2} + 2\,{\beta^2}\,{k_{y}^2} -
     {\omega^2} \right) \\ \mp 2\,\rho\,{\beta^2}\,k_{x}\,{k}_{z,P}(k_x,k_y) \\
  \mp 2\,\rho\,{\beta^2}\,k_{y}\,{k}_{z,P}(k_x,k_y) \end{array}\right]
  \epsilon_{P}
\end{equation}
with the  normalization factor ${\epsilon}_{P}$

\item the $\left|\pm,S, k_{x},k_{y}  \right\rangle$ eigen-vectors
are the   ${\stackrel{\mbox{\small downgoing}}{{}_{\mbox{\small
upgoing}}}}$  quasi-$SV-$waves
\begin{equation} \label{SVeigen}
\left|\pm,S, k_{x},k_{y}  \right\rangle=\left[\begin{array}{c}
i\,k_{x}\\
\mp i\,{k}_{z,S}(k_x,k_y)\\0\\
\mp 2\,\rho\,{\beta^2}\,k_{x}\,{k}_{z,S}(k_x,k_y) \\
  \rho \,\left({\omega^2} -2\,{\beta^2}\,{k_{x}^2} - {\beta^2}\,{k_{y}^2}  \right) \\
  -\rho\,{\beta^2}\,k_{x}\,k_{y} \end{array}\right]
  \epsilon_{S}
\end{equation}
with the  normalization factor ${\epsilon}_{S}$

\item the $\left|\pm,H, k_{x},k_{y}  \right\rangle$ eigen-vectors
are the   ${\stackrel{\mbox{\small downgoing}}{{}_{\mbox{\small
upgoing}}}}$  quasi-$SH-$waves
\begin{equation} \label{SHeigen}
\left|\pm,H, k_{x},k_{y}  \right\rangle=\left[\begin{array}{c}
\mp k_{y}\,{k}_{z,H}(k_x,k_y) \\-k_{x}\,k_{y} \\
  \hat{K}_{H}^2\\
2\,i\,k_{y}\,\rho \,\left({\beta^2}\,{k_{x}^2} +
     {\beta^2}\,{k_{y}^2} - {\omega^2} \right) \\
  \mp 2\,i\,\rho\,{\beta^2}\,k_{x}\,k_{y}\,{k}_{z,H}(k_x,k_y) \\
  \pm i\,{k}_{z,H}(k_x,k_y)\,\rho \,\left(  {\omega^2} -2\,{\beta^2}\,{k_{y}^2} \right) \end{array}\right]
  \epsilon_{H}
\end{equation}
with the  normalization factor ${\epsilon}_{H}$
\end{itemize}
We have called the eigen-vectors in (\ref{SVeigen}) and (\ref{SHeigen})
{\em quasi-\/}$SV$ and {\em quasi-\/}$SH$ waves respectively, because the
corresponding vector fields are {\em solenoidal\/} and are {\em in the
$(z,x)$ plane\/}
and {\em out of the $(z,x)$ plane\/} respectively, when $k_{y}=0$. However, our
choice of the  shear wavefield decomposition into quasi-inplane and quasi-out
of plane components
differs from the choice commonly used in seismology [e.g.
{\em Prange \& Toks\"oz}~(1990)]. Unlike the conventional decomposition,  the basis vectors
(\ref{SVeigen}) and (\ref{SHeigen}) are {\em non-degenerate\/} when
${k}_{x}=0$ and ${k}_{y}=0$.
The normalization factors ${\epsilon}_{P},{\epsilon}_{S} $ and
${\epsilon}_{H}$
in Eqns.(\ref{Peigen})--(\ref{SHeigen}) are chosen so that ${\bmi
D}_{z}^{-1}({k}_{x},{k}_{y})$
takes a simple form.

Using the standard result that eigenspaces of ${\bmi  A}^{T}({k}_{x},{k}_{y})$ annihilate
eigenspaces of ${\bmi  A}({k}_{x},{k}_{y})$ (and vice versa), we can avoid constructing  ${\bmi  D}_{z}^{-1}({k}_{x},{k}_{y})$ in (\ref{specA}) by explicit
matrix inversion, since we use the {\em Hamiltonian symmetry\/}
(\ref{ATdef}), to relate ${\bmi  D}_{z}^{-1}({k}_{x},{k}_{y})$ to the eigen matrix of the
{\em transposed\/} or {\em
dual\/} system matrix ${\bmi  A}^{T}({k}_{x},{k}_{y})$ defined as
\[
{\bmi  A}^{T}({k}_{x},{k}_{y})\equiv {\bmi  J}_6 {\bmi  A}(-{k}_{x},-{k}_{y}) {\bmi  J}_6
\]
To begin, it follows from (\ref{eigDef}) and (\ref{kzcDef}) that ${\bmi
\Lambda}(-{k}_{x},-{k}_{y})={\bmi \Lambda} ({k}_{x},{k}_{y})$,
so that
\begin{eqnarray*}
{\bmi  D}_{z}(-{k}_{x},-{k}_{y}){\bmi \Lambda} ({k}_{x},{k}_{y})&=
{\bmi  A}(-{k}_{x},-{k}_{y}){\bmi  D}_{z}(-{k}_{x},-{k}_{y})\\
&={\bmi  J}_{6}{\bmi
A}^{T}({k}_{x},{k}_{y}){\bmi  J}_{6}{\bmi  D}_{z}(-{k}_{x},-{k}_{y})
\end{eqnarray*}
which implies upon rearrangement that
\[
{\bmi  A}^{T}({k}_{x},{k}_{y}){\bmi  J}_{6}{\bmi  D}_{z}(-{k}_{x},-{k}_{y})=-{\bmi  J}_{6}
{\bmi  D}_{z}(-{k}_{x},-{k}_{y}){\bmi
\Lambda} ({k}_{x},{k}_{y}).
\]
Consequently, ${\bmi  J}_{6}{\bmi  D}_{z}(-{k}_{x},-{k}_{y})$ is an eigenbasis of the dual
operator ${\bmi
A}^{T}({k}_{x},{k}_{y})$, with $-{\bmi  \Lambda} ({k}_{x},{k}_{y})$ the corresponding eigenvalue array.
From direct evaluation it follows that
\begin{equation} \label{dinv1}
{\left({{\bmi  J}_{6}{\bmi  D}_{z}(-{k}_{x},-{k}_{y})}\right)}^{T}{\bmi
D}_{z}({k}_{x},{k}_{y})= i\left[{\begin{array}{cc}{\bf 0}_{3}&{\bmi V}\\
-{\bmi V}&{\bf 0}_{3}\end{array}}\right]
\end{equation}
where
\[
{\bmi V}=
\mbox{diag}\left[\begin{array}{ccc}
2 \rho {\omega }^{  2}{  k}_{  z, P }({k}_{x}){\epsilon
}_{P}^{  2},\,&
2 \mu {  k}_{  z, S }({k}_{x}) \hat{K}_{S }^{2}{\epsilon }_{S}^{
2},\,&
2 \rho {\omega }^{  2}{  k}_{  z, H }({k}_{x}) \hat{K}_{H }^{2}{\epsilon }_{H}^{  2}
\end{array}\right]
\]
and  $\hat{K}_{c}, c\in\{P,S,H\}$ was given by (\ref{Kdef}),
 so that if we choose the normalization factors ${\epsilon }_{P }$,
${\epsilon }_{S}$ and ${\epsilon }_{H}$ in Eqns.(\ref{Peigen})--(\ref{SHeigen})
according to
 \begin{equation} \label{epsdef}
{\epsilon }_{P }= \frac{1}{\sqrt{2 \rho {\omega }^{  2}{  k}_{
z, P }}}, \quad
{\epsilon }_{S }= \frac{\omega}{\beta\,\hat{K}_{S }\,\sqrt{2 \rho {\omega }^{  2}  {  k}_{
z, S }}}
, \quad \mbox{and} \quad
{\epsilon }_{H }= \frac{1}{\hat{K}_{H }\, \sqrt{2 \rho {\omega }^{  2}  {  k}_{
z, H }}}
\end{equation}
 then (\ref{dinv1}) simplifies to
\begin{equation} \label{dinv2}
{\left({{\bmi  J}_{6}{\bmi  D}_{z}(-{k}_{x},-{k}_{y})}\right)}^{T}{\bmi  D}_{z}({k}_{x},{k}_{y})
=i {\bmi  J}_{6}
\end{equation}
so that upon rearrangement of (\ref{dinv2}) we get
\begin{equation}	\label{D1def}
{\bmi  D}_{z}^{-1}({k}_{x},{k}_{y})=-i\	{\bmi  J}_6{\bmi  D}_{z}^{T}(-{k}_{x},-{k}_{y}){\bmi  J}_6
\end{equation}
Equation (\ref{D1def}) for the inverse is reminiscent of the {\em
symplectic\/} structures found in {\em Hamiltonian dynamical systems\/}
[e.g. {\em Goldstein}~(1980)].


We now transform the $z-$dependence of  (\ref{FLdef}) using
the {\em  Fourier transform\/}
\[
{\cal F}_{z}\,
[\cdot] ({k}_{z}) \equiv \frac{1}{\sqrt{2 \pi}}\int_{-\infty }^{ \infty }
dz\ \exp \left({-i\ {k}_{z}\ z}\right)\, [\cdot] (z) ,\]
 to get
\begin{equation}	\label{FzLdef}
	\tilde{\bmi  L}({\bmi k}) \,\rf{\bmi  G}({\bmi k};{\bmi x}^{\prime};k_y) \equiv \left[i
{\bmi  I}_6 {k}_{z} -{\bmi
A}({k}_{x},{k}_{y})\,
	\right]\,\rf{\bmi  G}({\bmi k};{\bmi x}^{\prime};k_y)={\bmi  I}_{6}
	\frac{\exp({-i {\bmi k}\cdot {\bmi x}^{\prime}})}{{2 \pi}},
\end{equation}
where ${\bmi x}^{\prime}=(z^{\prime},x^{\prime})$ and  ${\bmi
k}=(k_{z},k_{x})$, so that
\begin{equation}	\label{FzLdef1}
	\rf{\bmi  G}({\bmi k};{\bmi x}^{\prime};k_y) \equiv
	\left[i {\bmi  I}_6 {k}_{z} -{\bmi  A}({k}_{x},{k}_{y}) \right]^{-1}
	\frac{\exp({-i {\bmi k}\cdot {\bmi x}^{\prime}})}{{2 \pi}}.
\end{equation}
Using the spectral representation (\ref{specA}) in  (\ref{FzLdef1}),
we get
\begin{align} \label{InvLDef0}
\rf{\bmi  G}({\bmi k};{\bmi x}^{\prime};k_y)
&={\bmi  D}_{z}({k}_{x},{k}_{y})
{\left[i\, {\bmi  I}_{6}\, {k}_{z}\, -\, i\,   {\bmi \Lambda} \right]}^{-1}
{\bmi  D}_{z}^{-1}({k}_{x},{k}_{y}) \frac{\exp({-i {\bmi k}\cdot {\bmi x}^{\prime}})}{2 \pi}
\nonumber\\
&={\bmi  D}_{z}({k}_{x},{k}_{y})
{\left[{\bmi  J}_{6}\, {k}_{z}\, - \,{\bmi  J}_{6}\, {\bmi \Lambda}
\right]}^{-1}
{\bmi  D}_{z}^{T}(-{k}_{x},-{k}_{y})\,{\bmi  J}_6
\frac{\exp(-i {\bmi k}\cdot {\bmi x}^{\prime})}{2 \pi}\nonumber\\
&={\bmi  D}_{z}({k}_{x},{k}_{y})
\left[{\begin{array}{cc}{\bf 0}_{3}&-{\bmi W}^{+}\\
{\bmi W}^{-}&{\bf 0}_{3}\end{array}}\right]
{\bmi
D}_{z}^{T}(-{k}_{x},-{k}_{y})
\,{\bmi  J}_6 \frac{\exp(-i {\bmi k}\cdot {\bmi x}^{\prime})}{{2 \pi}}
\end{align}
where
\[
{\bmi W}^{\pm}=\mbox{diag}
\left[\begin{array}{ccc}
\frac{1}{({k}_{z}\mp{k}_{z,P })},\,&\frac{1}{({k}_{z}\mp{k}_{z,S
})},\,&\frac{1}{({k}_{z}\mp{k}_{z,S })}\end{array}\right]
\]


Applying the inverse Fourier transform
${\CB F}^{-1}\, [\cdot]({\bmi x}) \equiv \frac{1}{{2 \pi}}
\int_{R^2 } {\rm d}{\bmi k} \exp \left(i {\bmi k}\cdot {\bmi x}\right)\,
[\cdot] ({\bmi k})$,
it follows from (\ref{InvLDef0}),
that $\rf{\bmi  G}({\bmi x};{\bmi x}^{\prime};k_y)$
may be written
\begin{align} \label{InvLDef}
\rf{\bmi  G}({\bmi x};{\bmi x}^{\prime};k_y)&={\CB F}^{-1}\, \left[\rf{\bmi  G}({\bmi
k};{\bmi x}^{\prime};k_y) \right]({\bmi x}) =\frac{1}{4 \pi^2}
\int_{R^2 } {\rm d}{\bmi k} \exp \left[i {\bmi k}\cdot ({\bmi x}-{\bmi
x}^{\prime})\right]\nonumber\\
&\quad\times {\bmi  D}_{z}({k}_{x},{k}_{y})
\left[{\begin{array}{cc}{\bf 0}_{3}&-{\bmi W}^{+}\\
{\bmi W}^{-}&{\bf 0}_{3}\end{array}}\right]{\bmi  D}_{z}^{T}(-{k}_{x},-{k}_{y})
\,{\bmi  J}_6
\end{align}

\subsubsection*{3.1.1\quad The Classical Up/Down, Plane Wave Green's Tensor}
We now evaluate the $k_{z}$ integral in (\ref{InvLDef}) using the calculus of
residues [e.g. {\em Aki \& Richards\/}~(1980)], which is possible provided
the $z$ component of  $({\bmi x}-{\bmi x}^{\prime})$ in  (\ref{InvLDef})
is non-zero, i.e. $({z}-{z}^{\prime})\not=0 $.

We  close the
$k_z$ integration contour
in the ${\stackrel{{\mbox{{\small upper }}}}{{}_{\mbox{{\small lower }}}}}$ half of the
complex ${k}_{z}-$plane when
$(z-{z}^{\prime})\gtlt0$, in so doing  we pick up the
\[
\Re(\partial \omega/\partial { k}_{z,c})\gtlt 0
\]
residue contributions, associated with
${\stackrel{\mbox{\small downgoing}}{{}_{\mbox{\small upgoing}}}}$
group velocity, to get
%%% FIX: use multline to prevent the long equation from overflowing the page
\begin{align} \label{Gzres}
\rf{\bmi  G}({\bmi x};{\bmi x}^{\prime};k_y)&=\frac{- i}{2 \pi}
\int_{R} {\rm d}{ k}_{x} \exp \left[i { k}_{x} ({ x}-{
x}^{\prime})\right] \nonumber\\
&\quad\times\, {\bmi  D}_{z}({k}_{x},{k}_{y})
{\bmi Q}^{z}(z;z^{\prime}){\bmi  D}_{z}^{T}(-{k}_{x},-{k}_{y})
\,{\bmi  J}_6
\end{align}
where
${\bmi Q}^{z}(z;z^{\prime})$ in (\ref{Gzres}),
is defined  as
\begin{equation} \label{Qzdef}
{\bmi Q}^{z}(z;z^{\prime})=
\left[{\begin{array}{cc}{\bf 0}_{3}&H(z-{z}^{\prime}){\bmi E}(k_x;z;{z}^{\prime})\\
H(z^{\prime}-{z}){\bmi E}(k_x;z;{z}^{\prime})&{\bf 0}_{3}\end{array}}\right]
\end{equation}
with
\begin{equation} \label{Ek1def}
{\bmi E}(k_x;z;{z}^{\prime})=\left[{\begin{array}{ccc}
{\rm e}^{i{ k}_{z,P}(k_{x})|z-{z}^{\prime}|}&{ 0}&{ 0}\\
{ 0}&{\rm e}^{
i{ k}_{z,S}(k_{x})|z-{z}^{\prime}|}&{ 0}\\ { 0}&{ 0}&{\rm e}^{
i{ k}_{z,H}(k_{x})|z-{z}^{\prime}|}\end{array}}\right] ,
\end{equation}
and $H(\cdot)$ is the {\em Heaviside step function\/}.
The imposition of the {\em radiation boundary conditions\/} employed here,
follows {\em Lighthill}~(1960). Alternatively,  we may
 introduce {\em attenuation\/} to move the
poles off the integration contour [e.g. {\em Aki \& Richards\/}~(1980)].

\newpage

\subsection*{3.2\quad Discrete Plane Wave Representations of the Free Space Elastodynamic Green's Tensor}

\subsubsection*{3.2.1\quad The $z$ Splitting}

Basic results concerning the discrete Fourier (DFT) approximation to the inverse
Fourier transform over horizontal wavenumber of a scalar function,
 are summarized in {\sc Box} $\bf 3.2$ and a simple extension of the
 basic theory to incorporate Filon quadrature, while exploiting the speed
 of the FFT algorithm is presented in {\sc Box} $\bf 3.3$.
We will employ these results to develop a DFT approximation to (\ref{Gzres}).
We assume the $N_k=2 M$ element {\em Fourier lattice\/}
\begin{equation} \label{pspec}
\left\{
 {k}_{j} =j {\mit \Delta} k_{x} \ |\  j=-M,\cdots ,M-1
\right\}
\end{equation}
where ${\mit \Delta} k_{x} =2 \pi  j/L_{x}$ and   $L_{x}$ defines the {\em fundamental spatial period\/}
of the representation.
We may then approximate (\ref{Gzres}) by the {\em discrete
plane wave sum\/}
\begin{align} \label{Gdef5}
\rf{\bmi G}\!{}^{z}({\bmi x};{\bmi x}^{\prime};k_y)&=\frac{- i\,{\mit \Delta} {k}_{x}}{2 \pi}
\sum\limits_{{j=-M}}^{j=M-1}
\nonumber\\
&\quad\times
\exp \left[i {k}_{j} ({ x}-{
x}^{\prime})\right] {\bmi  D}_{z}({k}_{j},{k}_{y})
{\bmi Q}^{z}({k}_{j};z;z^{\prime}){\bmi  D}_{z}^{T}(-{k}_{j},-{k}_{y})
\,{\bmi  J}_6
\end{align}
The discrete representation (\ref{Gdef5}) is equivalent to assuming that the
transverse dependence of the wavefield is both {\em  bandlimited\/}
(since $N_k$ is
finite) and periodically repeated (since ${\mit \Delta} {k}_{x}$ is finite).


The next step will be to rewrite the discrete sum (\ref{Gdef5}) as an
equivalent matrix operation. To this end, we define the $N_k-$element
diagonal,
horizontal-wavevector array:
\begin{equation} \label{gKx}
{\bf  k}_{x} =\mbox{diag}\left[0,
{\mit \Delta} k_{x},\cdots, (M -1){\mit \Delta} k_{x},
-M {\mit \Delta} k_{x},\cdots,-{\mit \Delta} k_{x},\right];
\end{equation}
the $N_k-$element
diagonal,
vertical-wavevector array
\begin{equation} \label{gKz}
{\bf  k}_{z,c} =\mbox{diag}\left[{ k}_{z,c}(0),{ k}_{z,c}({\mit \Delta} k_{x}),\cdots,
{ k}_{z,c}((M-1){\mit \Delta} k_{x}),
{ k}_{z,c}(-M{\mit \Delta} k_{x}),
\cdots, { k}_{z,c}(-{\mit \Delta} k_{x})\right], \quad c\in\{P,S,H\}
\end{equation}
and the $N_k-$element {\em Fourier vector\/}
\begin{equation} \label{Fdef}
{\bmi F}(x;{\bf k}_{x})=
\left[\begin{array}{ccc} 1&, \cdots, &1 \end{array}
\right]\cdot{\rm e}^{i\,
{\bf k}_{x}\,x}.
\end{equation}
It is important to note that ${\bmi F}(x;{\bf k}_{x})$ as defined in
(\ref{Fdef}) is an $N_k-$ row vector, even though ${\bf k}_{x}$ is an
$N_k-$diagonal matrix.


Using the newly defined arrays (\ref{gKx})--(\ref{Fdef}),
we can  rewrite the {\em discrete plane wave\/}
approximation (\ref{Gdef5}) to the $6 \times 6$ Green's tensor in the equivalent matrix
form
\begin{align} \label{superG}
\rf{\bmi G}\!{}^{z}({\bmi x};{\bmi x}^{\prime};k_y)
&=- i\,\left({\bmi I}_6\otimes
{\bmi F}(x;{\bf k}_{x})  \right) {\bmi  D}_{z}({\bf k}_{x},{k}_{y})\nonumber\\
&\quad\times
{\bmi Q}^{z}({\bf k}_{x};z;z^{\prime}){\bmi  D}_{z}^{T}(-{\bf k}_{x},-{k}_{y})
\,\left({\bmi I}_6\otimes {\bmi F}^{T}(x^{\prime};-{\bf k}_{x})  \right)\,{\bmi  J}_6
\end{align}
where:\newline
the {\em Kronecker product\/}
$\,\left({\bmi I}_6\otimes {\bmi F}(x^{\prime};{\bf k}_{x})  \right)$ is
a $  6 \times 6 N_k$ array,
${\bmi F}^{H}(x;{\bf k}_{x})$ is the {\em Hermitian transpose\/} of the
Fourier vector ${\bmi F}(x;{\bf k}_{x})$ defined in (\ref{Fdef})
(so that ${\bmi F}^{H}(x;{\bf k}_{x})\equiv {\bmi F}^{T}(x;-{\bf
k}_{x})$);\newline
 the $(6N_k)\times (6N_k)$ matrix analogue of (\ref{Ddef}) is defined
as
\begin{equation} \label{Dsuperdef}
{\bmi D}_{z}({\bf k}_{x},k_{y})=\left[\begin{array}{cccccc}
\left|+,P, {\bf k}_{x},k_{y}  \right\rangle&\left|+,S, {\bf k}_{x},k_{y}
\right\rangle&\left|+,H, {\bf k}_{x},k_{y}  \right\rangle&
\left|-,P, {\bf k}_{x},k_{y}  \right\rangle&\left|-,S, {\bf k}_{x},k_{y}
\right\rangle&\left|-,H, {\bf k}_{x},k_{y}  \right\rangle
\end{array}\right]
\end{equation}
in terms of
\begin{itemize}
\item the $6 N_k \times N_k$ $P$-wave partition
\begin{equation} \label{PeigenSuper}
\left|\pm,P, {\bf k}_x,k_{y}  \right\rangle=
\sqrt{\frac{{\mit \Delta} k_{x}}{{4 \pi \rho {\omega
}^{  2}{\bf  k}_{
z, P }}}}\,\left[\begin{array}{c} \pm i\,{\bf k}_{z,P}\\i\,{\bf k}_{x}\\
i\,k_{y}\,{\bmi I}\\
  \rho \,\left( 2\,{\beta^2}\,{{\bf k}_{x}^2} + ( 2\,{\beta^2}\,{k_{y}^2} -
     {\omega^2}){\bmi I} \right) \\ \mp 2\,\rho\,{\beta^2}\,{\bf k}_{x}\,{\bf k}_{z,P} \\
  \mp 2\,\rho\,{\beta^2}\,k_{y}\,{\bf k}_{z,P} \end{array}\right]
\end{equation}

\item  the $6 N_k \times N_k$ $SV$-wave partition
\begin{equation} \label{SVeigenSuper}
\left|\pm,S, {\bf k}_x,k_{y}  \right\rangle=\frac{\omega }{\beta \hat{K}_{S}}\,
\sqrt{\frac{{\mit \Delta} k_{x}}{{4 \pi \rho {\omega
}^{  2}{\bf  k}_{
z, S }}}}\,\left[\begin{array}{c}
i\,{\bf k}_{x}\\
\mp i\,{\bf k}_{z,S}\\ {\bf 0} \\
\mp 2\,\rho\,{\beta^2}\,{\bf k}_{x}\,{\bf k}_{z,S} \\
  \rho \,\left(({\omega^2} - {\beta^2}\,{k_{y}^2}) {\bmi I}-2\,{\beta^2}\,{{\bf k}_{x}^2}   \right) \\
  -\rho\,{\beta^2}\,k_{y}\,{\bf k}_{x} \end{array}\right]
\end{equation}

\item the $6 N_k \times N_k$ $SH$-wave partition
\begin{equation} \label{SHeigenSuper}
\left|\pm,H, {\bf k}_x, k_{y}  \right\rangle=\frac{1 }{\hat{K}_{H}}\,
\sqrt{\frac{{\mit \Delta} k_{x}}{{4 \pi \rho {\omega
}^{  2}{\bf  k}_{
z, H }}}}\,\left[\begin{array}{c}
\mp k_{y}\,{\bf k}_{z,H} \\-k_{y}\,{\bf k}_{x} \\
  \hat{K}_{H}^2 {\bmi I}\\
-2\,i\,\rho \,{\beta^2}\,k_{y}\,{{\bf k}_{z,H}^2}\\
  \mp 2\,i\,\rho\,{\beta^2}\,k_{y}\,{\bf k}_{x}\,{\bf k}_{z,H} \\
  \pm i\,\rho\,\left(  {\omega^2} -2\,{\beta^2}\,{k_{y}^2}
  \right)\,{\bf k}_{z,H}  \end{array}\right]
\end{equation}

\end{itemize}
The $(6N_k)\times (6N_k)$ {\em plane wave propagator\/}
${\bmi Q}^{z}({\bf k}_{x};z;z^{\prime})$ in (\ref{superG}),
is defined  as
\begin{equation} \label{Qdef}
{\bmi Q}^{z}({\bf k}_x;z;z^{\prime})=
\left[{\begin{array}{cc}{\bf 0}_{3N_k}&H(z-{z}^{\prime}){\bmi E}({\bf k}_x;z;{z}^{\prime})\\
H(z^{\prime}-{z}){\bmi E}({\bf k}_x;z;{z}^{\prime})&{\bf 0}_{3N_k}\end{array}}\right]
\end{equation}
where the $(3N_k)\times (3N_k)$ {\em vertical phase shift array\/} is
defined as
\[
{\bmi E}({\bf k}_x;z;{z}^{\prime})=\left[{\begin{array}{ccc}
{\rm e}^{i{\bf k}_{z,P}|z-{z}^{\prime}|}&{\bf 0}_{N_k}&{\bf 0}_{N_k}\\
{\bf 0}_{N_k}&{\rm e}^{
i{\bf k}_{z,S}|z-{z}^{\prime}|}&{\bf 0}_{N_k}\\ {\bf 0}_{N_k}&{\bf 0}_{N_k}&{\rm e}^{
i{\bf k}_{z,H}|z-{z}^{\prime}|}\end{array}}\right] ,
\]
in terms of the diagonal wavevector array ${\bf k}_{z,c}, c\in\{P,S,H
\}$ of (\ref{gKz}).

It will be convenient in later applications,  to split (\ref{superG}) into the product
\[
\rf{\bmi G}\!{}^{z}({\bmi x};{\bmi x}^{\prime};k_y)=
\rf{\bmi P}\!{}^{z}({\bmi x};{\bmi x}^{\prime};k_y)\cdot{\bmi J}_6
\]
where ${\bmi J}_6$ is the $6\times 6$ {\em canonically symplectic\/} array and
\begin{align} \label{superP}
\rf{\bmi P}\!{}^{z}({\bmi x};{\bmi x}^{\prime};k_y)
&=- i\,\left({\bmi I}_6\otimes
{\bmi F}(x;{\bf k}_{x})  \right) {\bmi  D}_{z}({\bf k}_{x},{k}_{y})\nonumber\\
&\quad\times
{\bmi Q}^{z}({\bf k}_{x};z;z^{\prime}){\bmi  D}_{z}^{T}(-{\bf k}_{x},-{k}_{y})
\,\left({\bmi I}_6\otimes {\bmi F}^{T}(x^{\prime};-{\bf k}_{x})  \right).
\end{align}
might be called the {\em propagator\/}, for the want of a better name. The
propagator is an entity in itself for most numerical applications
considered here.

The representation (\ref{superG}) (or equivalently, (\ref{superP})) is readily extensible
to the plane stratified problem.

\subsubsection*{3.2.2\quad The $x$ Splitting}


Recalling the definition (\ref{Ddef}) for the columns of the eigen-matrix
${\bmi  D}_{z}$, (\ref{InvLDef}) may be rewritten as
\begin{align}\label{grep2}
\rf{\bmi  G}({\bmi  x};{{\bmi  x}}^{\prime};k_y)
&=\sum\limits_{c\in \{P,S,H\}}
{\frac{1}{4{ \pi }^{2}}}\int_{{R}^{2}} {\rm d}{\bmi k} \exp \left[{i {\bmi k}  \cdot
  ({\bmi x}-{\bmi x}^{\prime})}\right]\nonumber\\
&\quad\times\Biggl[\frac{
  \left|-,c, k_{x},k_{y}  \right\rangle\otimes \left|+,c, -k_{x},-k_{y}  \right\rangle
  }{{k}_{z}+{k}_{z, c }({k}_{x})}\nonumber\\
&\qquad\qquad{}-\frac{
\left|+,c, k_{x},k_{y}  \right\rangle\otimes \left|-,c, -k_{x},-k_{y}  \right\rangle
  }{{k}_{z}-{k}_{z, c }({k}_{x})}\Biggr]\,{\bmi J}_{6}
\nonumber\\
&=
\rf{\bmi  G}_{P}({\bmi  x};{{\bmi  x}}^{\prime};k_y)+
\rf{\bmi  G}_{S}({\bmi  x};{{\bmi  x}}^{\prime};k_y)+
\rf{\bmi  G}_{H}({\bmi  x};{{\bmi  x}}^{\prime};k_y)
\end{align}
where  $\otimes$ is the {\em tensor product\/}.
We will ultimately evaluate the integrals in (\ref{grep2}) using the
residues at the simple poles, where
\begin{equation} \label{dispersion}
{k}^2_{z}+{k}^2_{x}+{k}^2_{y}=\frac{\omega^2}{c^2},\quad \mbox{for}\quad c\in\{P,S,H\}.
\end{equation}
Hence,  (\ref{grep2}) may be replaced by an expression that is
{\em equivalent at these poles\/}. Introducing the partial differential
operators
\begin{align} \label{Labdef}
\begin{array}{lll}
{\bmi L}_{P}=&\qquad{\bmi L}_{S}=&\qquad{\bmi L}_{H}=\\
\left[\begin{array}{c}
\partial_{z} \\
\partial_{x} \\
\partial_{y} \\
2 \mu \partial_{zz} - {\lambda \omega^{2}\over \alpha^2}\\
    2 \mu \partial_{xz} \\
   2 \mu \partial_{yz}
\end{array}\right]&\qquad
\frac{\omega }{\beta \hat{K}_{S}}\,
\left[\begin{array}{c}
\partial_{x} \\-\partial_{z} \\ 0 \\
2 \mu \partial_{xz}\\
   \mu (-\partial_{zz} + \partial_{xx}) \\
   \mu \partial_{xy}
\end{array}\right]&\qquad
\frac{1 }{\hat{K}_{H}}\,
\left[\begin{array}{c}
   \partial_{yz} \\
\partial_{xy} \\
{\omega^{2}\over \beta^2} + \partial_{yy} \\
 2 \mu \partial_{yzz}\\
    2 \mu \partial_{xyz} \\
   \mu ({\omega^{2}\over \beta^2}\partial_{z} + 2 \partial_{yyz})
\end{array}\right]
\end{array}
\end{align}
$\rf{\bmi G}_{c}({\bmi  x};{{\bmi  x}}^{\prime};k_y)$  in  (\ref{grep2}) may be replaced
with the {\em equivalent\/} expression
\begin{align}\label{grepc}
\rf{\bmi G}_{c}({\bmi  x};{{\bmi  x}}^{\prime};k_y)
&= {\frac{1}{8{ \pi }^{2} \rho \omega^2}}\int_{{R}^{2}} {\rm d}{\bmi k}\,
\frac{
{\bmi L}_{c} \otimes {\bmi  L}_{c}^{\prime}
\left(
\exp \left[{i {\bmi k}  \cdot  ({\bmi x}-{\bmi x}^{\prime})}\right]\right)
}{ {k}_{z, c }({k}_{x})\,[{k}_{z}+{k}_{z, c }({k}_{x})]} {\bmi J}_{6}\nonumber\\
&\quad -
{\frac{1}{8  { \pi }^{2}\rho \omega^2}}\int_{{R}^{2}} {\rm d}{\bmi k}\,
\frac{{\bmi L}_{c} \otimes {\bmi  L}_{c}^{\prime}\left(
\exp \left[{i {\bmi k}  \cdot  ({\bmi x}-{\bmi x}^{\prime})}\right]\right)
}{{k}_{z, c }({k}_{x})\, [{k}_{z}-{k}_{z, c }({k}_{x})]} {\bmi J}_{6},\nonumber\\
&\hspace{20em}
c\in\{ P, S , H\};\nonumber\\
&= \frac{1}{4{ \pi }^{2}\rho \omega^2}
\int_{{R}^{2}} {\rm d}{\bmi k}\, {\bmi L}_{c} \otimes {\bmi  L}_{c}^{\prime}
\frac{
\exp \left[{i {\bmi k}  \cdot  ({\bmi x}-{\bmi x}^{\prime})}\right]
}{ \hat{K}^{2}_{ c } - {k}^{2}_{z} - {k}^{2}_{x}} {\bmi J}_{6}
\end{align}
where we have used the algebraic relation (\ref{dispersion}) and
 the differential operator ${\bmi  L}_{c}^{\prime}$ operates on the
${\bmi  x}^{\prime}$ coordinate.


Assuming that $x \not = x^{\prime}$ in (\ref{grepc}) we may  evaluate the $k_x$ integral in
that equation using the calculus of residues to get
\begin{align} \label{grep3x}
\rf{\bmi G}_{c}({\bmi  x};{{\bmi  x}}^{\prime};k_y)&=
\frac{-i}{2{ \pi }\rho \omega^2}
\int_{{R}} {\rm d}{k}_z \,{\bmi L}_{c } \otimes {\bmi  L}_{c}^{\prime}\nonumber\\
&\quad\times\frac{\exp \left[{i {k}_z  ({z}-{z}^{\prime})}\right]
\exp \left[{i {k}_{x,c}({k}_z,{k}_y)  |{x}-{x}^{\prime}|}\right]
}{ 2\,{k}_{x,c}({k}_z,{k}_y)}{\bmi J}_{6},
\end{align}
with
\begin{equation}	\label{kxcDef}
	{k}_{x,c}({k}_{z},{k}_{y}) = \left\{
	\begin{array}{c}
	\sqrt {\hat{ K}_{c}^{2}({k}_{y})\ -\ {k}_{z}^{2}}\ \ \ \ \ \ |{k}_{x}|\ <\
	|\hat{K}_{c}({k}_{y})|\\
	i\sqrt {{k}_{z}^{ 2}\ -\ \hat{ K}_{c}^{2}({k}_{y})} \ \ \ \ |{k}_{x}|\ \ge   \
	|\hat{K}_{c}({k}_{y})|
\end{array}\right.
\end{equation}
and $\hat{ K}_{c}({k}_{y})$ was given by (\ref{Kdef}).
Using (\ref{grep3x}) in (\ref{grep2})
we get
\begin{align} \label{GxGRep}
\rf{\bmi G}_{c}({\bmi  x};{{\bmi  x}}^{\prime};k_y)&=
\frac{-i}{2{ \pi }\rho \omega^2}\,\sum\limits_{c\in \{P,S,H\}}\nonumber\\
&\quad\times
\int_{{R}} {\rm d}{k}_z \,{\bmi L}_{c } \otimes {\bmi  L}_{c}^{\prime}
\frac{\exp \left[{i {k}_z  ({z}-{z}^{\prime})}\right]
\exp \left[{i {k}_{x,c}({k}_z,{k}_y)  |{x}-{x}^{\prime}|}\right]
}{ 2\,{k}_{x,c}({k}_z,{k}_y)}{\bmi J}_{6}
\end{align}

Using the definition of the
partial differential operators ${\bmi L}_{c }$ from (\ref{Labdef}) in (\ref{grep2}) and
performing the relevant spatial differentiations  we get
\begin{align} \label{Gxres}
\rf{\bmi G}({\bmi x};{\bmi x}^{\prime};k_y)&=\frac{- i}{2 \pi}
\int_{R} {\rm d}{ k}_{z} \exp \left[i { k}_{z} ({ z}-{
z}^{\prime})\right]\nonumber\\
&\quad\times\, {\bmi  D}_{x}({k}_{z},{k}_{y})
{\bmi Q}^{x}(k_z;x;x^{\prime}){\bmi  D}_{x}^{T}(-{k}_{z},-{k}_{y})
\,{\bmi  J}_6
\end{align}
provided $x \not = x^{\prime}$, where:\newline
\begin{equation} \label{Dxdef}
{\bmi D}_{x}(k_{z},k_{y})=\left[\begin{array}{cccccc}
\overline{\left|+,P, k_{z},k_{y}  \right\rangle}&\overline{\left|+,S, k_{z},k_{y}
\right\rangle}&\overline{\left|+,H, k_{z},k_{y}  \right\rangle}&
\overline{\left|-,P, k_{z},k_{y}  \right\rangle}&\overline{\left|-,S, k_{z},k_{y}
\right\rangle}&\overline{\left|-,H, k_{z},k_{y}  \right\rangle}
\end{array}\right]
\end{equation}
where:
\begin{itemize}
\item the $\overline{\left|\pm,P, k_{z},k_{y}  \right\rangle}$ states
are the   ${\stackrel{\mbox{\small rightgoing}}{{}_{\mbox{\small
leftgoing}}}}$  $P-$waves
\begin{equation} \label{Pxeigen}
\overline{\left|\pm,P, k_{x,P},k_{y}  \right\rangle}=\frac{1}{\sqrt{2 \rho {\omega }^{  2}{  k}_{
x, P }}}\,\left[\begin{array}{c}
i\,{k}_{z}\\ \pm i\,k_{x,P}(k_x,k_y)\\i\,k_{y}\\
  \rho \,\left( 2\,{\beta^2}\,{k_{x,P}^2(k_x,k_y)} + 2\,{\beta^2}\,{k_{y}^2} -
     {\omega^2} \right) \\ \mp 2\,\rho\,{\beta^2}\,k_{x,P}(k_x,k_y)\,{k}_{z} \\
  - 2\,\rho\,{\beta^2}\,k_{y}\,{k}_{z} \end{array}\right]
\end{equation}

\item the $\overline{\left|\pm,S, k_{z},k_{y}  \right\rangle}$ states
are the   ${\stackrel{\mbox{\small rightgoing}}{{}_{\mbox{\small
leftgoing}}}}$  quasi-$SV-$waves
\begin{equation} \label{SVxeigen}
\overline{\left|\pm,S, k_{z},k_{y}  \right\rangle}=\frac{\omega}{\beta\,\hat{K}_{S }\,
\sqrt{2 \rho {\omega }^{  2}  {  k}_{
x, S }}}\,
\left[\begin{array}{c}
\pm i\,k_{x,S}(k_x,k_y)\\
- i\,{k}_{z}\\0\\
\mp 2\,\rho\,{\beta^2}\,k_{x,S}(k_x,k_y)\,{k}_{z} \\
  \rho \,\left({\omega^2} -2\,{\beta^2}\,{k_{x,S}^2(k_x,k_y)} - {\beta^2}\,{k_{y}^2}  \right) \\
  \mp \rho\,{\beta^2}\,k_{x,S}(k_x,k_y)\,k_{y} \end{array}\right]
\end{equation}

\item the $\overline{\left|\pm,H, k_{z},k_{y}  \right\rangle}$ states
are the   ${\stackrel{\mbox{\small rightgoing}}{{}_{\mbox{\small
leftgoing}}}}$  quasi-$SH-$waves
\begin{equation} \label{SHxeigen}
\overline{\left|\pm,H, k_{z},k_{y}  \right\rangle}=
\frac{1}{\hat{K}_{H }\, \sqrt{2 \rho {\omega }^{  2}  {  k}_{
x, H }}} \,\left[\begin{array}{c}
- k_{y}\,{k}_{z} \\\mp k_{x,H}(k_x,k_y)\,k_{y} \\
  \hat{K}_{H}^2\\
2\,i\,k_{y}\,\rho \,\left({\beta^2}\,{k_{x,H}^2(k_x,k_y)} +
     {\beta^2}\,{k_{y}^2} - {\omega^2} \right) \\
  \mp  2\,i\,\rho\,{\beta^2}\,k_{x,H}(k_x,k_y)\,k_{y}\,{k}_{z}\\
   i\,{k}_{z}\,\rho \,\left(  {\omega^2} -2\,{\beta^2}\,{k_{y}^2} \right) \end{array}\right]
\end{equation}
\end{itemize}
the {\em horizontal plane wave propagator\/} ${\bmi Q}^{x}(k_z;x;x^{\prime})$ in (\ref{Gxres}),
is defined  as
\begin{equation} \label{Qxdef}
{\bmi Q}^{x}(k_z;x;x^{\prime})=
\left[{\begin{array}{cc}{\bf 0}_{3}&H(x-{x}^{\prime}){\bmi E}(k_z;x;{x}^{\prime})\\
H(x^{\prime}-{x}){\bmi E}(k_z;x;{x}^{\prime})&{\bf 0}_{3}\end{array}}\right]
\end{equation}
with the {\em horizontal phase shift\/} array
\[
{\bmi E}(k_z;x;{x}^{\prime})=\left[{\begin{array}{ccc}
{\rm e}^{i{ k}_{x,P}(k_{z})|x-{x}^{\prime}|}&{ 0}&{ 0}\\
{ 0}&{\rm e}^{
i{ k}_{x,S}(k_{z})|x-{x}^{\prime}|}&{ 0}\\ { 0}&{ 0}&{\rm e}^{
i{ k}_{x,H}(k_{z})|x-{x}^{\prime}|}\end{array}}\right] ,
\]
and $H(\cdot)$ is the {\em Heaviside step function\/}.


For computational purposes, we must employ a {\em discrete plane wave\/}
or {\em Fourier lattice\/} version of the wavenumber integral
(\ref{Gxres}). The development closely parallels that of \S 3.2.1.
We will now derive a {\em discrete plane wave\/} version of
(\ref{Gxres}), and we note that we
will employ (\ref{Gxres}) only in the special case where $z=z^{\prime}$
so that $\exp \left[i \overline{k}_{j} ({ z}-{
z}^{\prime})\right]\equiv 1$ and consequently
\begin{align} \label{Gxres0}
\rf{\bmi G}({ x};{ x}^{\prime};k_y)&=\frac{- i}{2 \pi}
\int_{R} {\rm d}{ k}_{z}\,  {\bmi  D}_{x}({k}_{z},{k}_{y})\nonumber\\
&\quad\times\,{\bmi Q}^{x}(k_z;x;x^{\prime}){\bmi  D}_{x}^{T}(-{k}_{z},-{k}_{y})
\,{\bmi  J}_6, \quad z=z^{\prime}, \ x\not=x^{\prime}
\end{align}
so there will be no {\em
aliasing\/} or {\em wraparound\/} in the $z-$coordinate. This observation
is crucial to the success in the intended application of $\rf{\bmi G}({\bmi
x};{\bmi x}^{\prime};k_y)$ from (\ref{Gxres0}) to problems with a
$z-$stratification in material properties.


Given the $\overline{N}_k=2 \overline{M}+1$ element {\em Fourier lattice\/}
\begin{equation} \label{pxspec}
\left\{
 \overline{k}_{j} =j {\mit \Delta} k_{z} \ |\  j=-\overline{M},\cdots
 ,\overline{M}
\right\}
\end{equation}
where ${\mit \Delta} k_{z} =2 \pi  j/L_{z}$ and   $L_{z}$
controls the  {\em wavenumber sampling rate,\/} we may approximate
(\ref{Gxres0}) by the {\em discrete
plane wave sum\/}
\begin{align} \label{Gxdef5}
\rf{\bmi G}\!{}^{x}({ x};{ x}^{\prime};k_y)&=\frac{- i\,{\mit \Delta} {k}_{z}}{2 \pi}
\sum\limits_{{j=-\overline{M}}}^{j=\overline{M}}
{\bmi  D}_{x}(\overline{k}_{j},{k}_{y})\nonumber\\
&\quad\times\,{\bmi Q}^{x}(\overline{k}_{j};x;x^{\prime}){\bmi  D}_{x}^{T}(-\overline{k}_{j},-{k}_{y})
\,{\bmi  J}_6
\end{align}
The discrete representation (\ref{Gxdef5}) is equivalent to assuming that the
$z-$dependence of the wavefield is {\em  bandlimited\/} (since $\overline{M}$ is
finite), which is certainly not the case in the limit $x \leftarrow
x^{\prime}$.


The next step will be to rewrite the discrete sum (\ref{Gxdef5}) as an
equivalent matrix operation. To this end, we define the diagonal wave
vector arrays:
\begin{equation} \label{gbarKz}
{\bf  k}_{z} =\mbox{diag}\left[{-\overline{M} {\mit \Delta} k_{z},\cdots,-{\mit
\Delta} k_{z},0,
{\mit \Delta} k_{z},\cdots, \overline{M} {\mit \Delta} k_{z}}\right];
\end{equation}
\begin{equation} \label{gbarKx}
{\bf  k}_{x,c} =\mbox{diag}\left[{{ k}_{x,c}(-\overline{M} {\mit \Delta} k_{z}),\cdots,
{ k}_{x,c}(-{\mit \Delta} k_{z}),
{ k}_{x,c}(0),
{ k}_{z,c}({\mit \Delta} k_{z}),\cdots, { k}_{x,c}(\overline{M}{\mit \Delta} k_{z})}\right]
\end{equation}
and the $\overline{N}_k-$element {\em projection \/} vector
\begin{equation} \label{Fbardef}
\overline{\bmi F}(0)=\left[\begin{array}{ccccc}
1& \cdots &
1& \cdots &
1\end{array}\right].
\end{equation}


We can  rewrite the {\em discrete plane wave\/}
approximation (\ref{Gxdef5}) to the $6 \times 6$ Green's tensor in the equivalent matrix
form
\begin{align} \label{superGx}
\rf{\bmi G}\!{}^{x}({ x};{ x}^{\prime};k_y)
&=- i\,
\left({\bmi I}_6\otimes \overline{\bmi F}(0)  \right)
 {\bmi  D}_{x}({\bf k}_{z},{k}_{y})\nonumber\\
&\quad\times
{\bmi Q}^{x}({\bf k}_{z};x;x^{\prime}){\bmi  D}_{x}^{T}(-{\bf k}_{z},-{k}_{y})
\,\left({\bmi I}_6\otimes \overline{\bmi F}^{T}(0)  \right)\,{\bmi  J}_6
\end{align}
where:\newline
the $(6\bar{N}_k)\times (6\bar{N}_k)$ matrix analogue of (\ref{Dxdef}) is defined
as
\begin{equation} \label{Dxsuperdef}
{\bmi D}_{x}({\bf k}_{z},k_{y})=\left[\begin{array}{cccccc}
\overline{\left|+,P, {\bf k}_{z},k_{y}  \right\rangle}&\overline{\left|+,S, {\bf k}_{z},k_{y}
\right\rangle}&\overline{\left|+,H, {\bf k}_{z},k_{y}  \right\rangle}&
\overline{\left|-,P, {\bf k}_{z},k_{y}  \right\rangle}&\overline{\left|-,S, {\bf k}_{z},k_{y}
\right\rangle}&\overline{\left|-,H, {\bf k}_{z},k_{y}  \right\rangle}
\end{array}\right]
\end{equation}
in terms of
\begin{itemize}
\item the $6 \overline{N}_k \times \overline{N}_k$ $P$-wave partition
\begin{equation} \label{PxeigenSuper}
\overline{\left|\pm,P, {\bf k}_z,k_{y}  \right\rangle}=
\sqrt{\frac{{\mit \Delta} k_{z}}{{4 \pi \rho {\omega
}^{  2}{\bf  k}_{
x, P }}}}
\left[\begin{array}{c} i\,{\bf k}_{z}\\ \pm i\,{\bf k}_{x,P}\\
i\,k_{y}\,{\bmi I}\\
  \rho \,\left( 2\,{\beta^2}\,{{\bf k}_{x,P}^2} + ( 2\,{\beta^2}\,{k_{y}^2} -
     {\omega^2}){\bmi I} \right) \\ \mp 2\,\rho\,{\beta^2}\,{\bf k}_{x,P}\,{\bf k}_{z} \\
  - 2\,\rho\,{\beta^2}\,k_{y}\,{\bf k}_{z} \end{array}\right]
\end{equation}

\item  the $6 \overline{N}_k \times \overline{N}_k$ $SV$-wave
\begin{equation} \label{SVxeigenSuper}
\overline{\left|\pm,S, {\bf k}_z,k_{y}  \right\rangle}=
\frac{\omega }{\beta \hat{K}_{S}}\,\sqrt{\frac{{\mit \Delta} k_{z}}{{4 \pi \rho {\omega
}^{  2}{\bf  k}_{
x, S }}}}\,\left[\begin{array}{c}
\pm i\,{\bf k}_{x,S}\\
- i\,{\bf k}_{z}\\ {\bf 0} \\
\mp 2\,\rho\,{\beta^2}\,{\bf k}_{x,S}\,{\bf k}_{z} \\
  \rho \,\left(({\omega^2} - {\beta^2}\,{k_{y}^2}) {\bmi
  I}-2\,{\beta^2}\,{{\bf k}_{x,S}^2}   \right) \\
 \mp \rho\,{\beta^2}\,k_{y}\,{\bf k}_{x,S} \end{array}\right]
\end{equation}

\item the $6 \overline{N}_k \times \overline{N}_k$ $SH$-wave partition
\begin{equation} \label{SHxeigenSuper}
\overline{\left|\pm,H, {\bf k}_z, k_{y}  \right\rangle}=\frac{1 }{\hat{K}_{H}}\,
\sqrt{\frac{{\mit \Delta} k_{z}}{{4 \pi \rho {\omega
}^{  2}{\bf  k}_{
x, H }}}}\,\left[\begin{array}{c}
- k_{y}\,{\bf k}_{z} \\ \mp k_{y}\,{\bf k}_{x,H} \\
  \hat{K}_{H}^2 {\bmi I}\\
-2\,i\,\rho \,{\beta^2}\,k_{y}\,{{\bf k}_{z}^2}\\
  \mp 2\,i\,\rho\,{\beta^2}\,k_{y}\,{\bf k}_{x,H}\,{\bf k}_{z} \\
   i\,\rho\,\left(  {\omega^2} -2\,{\beta^2}\,{k_{y}^2}
  \right)\,{\bf k}_{z}  \end{array}\right]
\end{equation}

\end{itemize}
where  ${\bf  k}_{x}$ and ${\bf  k}_{z,c}$ were given in
(\ref{gbarKx}) and (\ref{gbarKz}) respectively.

The $(6\bar{N}_k)\times (6\bar{N}_k)$ {\em plane wave propagator\/}
${\bmi Q}^{x}({\bf k}_{z};x;x^{\prime})$ in (\ref{superGx}),
is defined  as
\begin{equation} \label{gQxdef}
{\bmi Q}^{x}({\bf k}_z;x;x^{\prime})=
\left[{\begin{array}{cc}{\bf 0}_{3\bar{N}_k}&H(x-{x}^{\prime}){\bmi E}({\bf
k}_z;x;{x}^{\prime})\\
H(x^{\prime}-{x}){\bmi E}({\bf k}_z;x;{x}^{\prime})&{\bf 0}_{3\bar{N}_k}\end{array}}\right]
\end{equation}
where the $(3\bar{N}_k)\times (3\bar{N}_k)$ {\em vertical phase shift array\/} is
defined as
\[
{\bmi E}({\bf k}_z;x;{x}^{\prime})=\left[{\begin{array}{ccc}
{\rm e}^{i{\bf k}_{z,P}|x-{x}^{\prime}|}&{\bf 0}_{\overline{N}_k}&{\bf 0}_{\overline{N}_k}\\
{\bf 0}_{\overline{N}_k}&{\rm e}^{
i{\bf k}_{z,S}|x-{x}^{\prime}|}&{\bf 0}_{\overline{N}_k}\\ {\bf 0}_{\overline{N}_k}&{\bf 0}_{\overline{N}_k}&
{\rm e}^{
i{\bf k}_{z,H}|x-{x}^{\prime}|}\end{array}}\right] ,
\]
in terms of the diagonal wavevector array ${\bf k}_{x,c}, c\in\{P,S,H
\}$ of (\ref{gbarKx}).

For computational efficiency, we split (\ref{superGx}) into
\[
\rf{\bmi G}\!{}^{x}({ x};{ x}^{\prime};k_y)=\rf{\bmi P}\!{}^{x}({ x};{ x}^{\prime};k_y){\bmi  J}_6
\]
where
\begin{align} \label{superPx}
\rf{\bmi P}\!{}^{x}({ x};{ x}^{\prime};k_y)
&=- i\,
\left({\bmi I}_6\otimes \overline{\bmi F}(0)  \right)
 {\bmi  D}_{x}({\bf k}_{z},{k}_{y})\nonumber\\
&\quad\times
{\bmi Q}^{x}({\bf k}_{z};x;x^{\prime}){\bmi  D}_{x}^{T}(-{\bf k}_{z},-{k}_{y})
\,\left({\bmi I}_6\otimes \overline{\bmi F}^{T}(0)  \right)
\end{align}


\end{document}
